
\documentclass{article}
\usepackage{colortbl}
\usepackage{makecell}
\usepackage{multirow}
\usepackage{supertabular}

\begin{document}

\newcounter{utterance}

\centering \large Interaction Transcript for game `eqbench', experiment `eq\_bench', episode 42 with reflexion{-}Qwen3.5{-}27B.
\vspace{24pt}

{ \footnotesize  \setcounter{utterance}{1}
\setlength{\tabcolsep}{0pt}
\begin{supertabular}{c@{$\;$}|p{.15\linewidth}@{}p{.15\linewidth}p{.15\linewidth}p{.15\linewidth}p{.15\linewidth}p{.15\linewidth}}
    \# & \multicolumn{2}{c}{Player} && \multicolumn{2}{c}{Game Master} \\
    \hline

    \theutterance \stepcounter{utterance}  
    & & & \multicolumn{4}{p{0.6\linewidth}}{
        \cellcolor[rgb]{0.9,0.9,0.9}{
            \makecell[{{p{\linewidth}}}]{
                \texttt{\tiny{[P1$\langle$GM]}}
                \texttt{You have attempted this same task before and did not succeed. These are the plans you wrote after those attempts.} \\
\texttt{Trial \#1: New plan: My previous attempt failed because I assigned a high score to "Inadequate" (8) and a low score to "Angry" (2), which likely contradicted the specific emotional nuance of the dialogue where James expresses frustration and defensiveness rather than just self{-}doubt. James's lines ("I can't believe we're doing this again," "It's... complicated," "haunted") suggest a mix of resentment and anxiety, while Rachel's teasing ("terrible at expressing your feelings") provokes a defensive reaction. Next time, I will re{-}evaluate the dialogue to prioritize "Angry" or "Frustrated" (mapped to the available options) as the primary emotion, likely increasing the score for "Angry" to reflect his defensive tone and decreasing "Inadequate" to reflect that his struggle is more about emotional blockage than a lack of self{-}worth. I will ensure the scores sum to a logical distribution where the dominant emotion is clearly identified based on the specific textual evidence of his defensiveness.} \\
\\ 
\texttt{Trial \#2: New plan: My previous attempt failed because I over{-}indexed on "Inadequate" (7) while underestimating the defensive and frustrated undertones in James's dialogue. While James admits to being "intimidated" and "hard" to open up, his opening line ("I can't believe we're doing this again") and his description of being "haunted" suggest a deeper sense of resentment and frustration rather than simple self{-}doubt. Rachel's teasing about him being "terrible at expressing feelings" likely triggers a defensive reaction that aligns more with "Angry" (or frustrated anger) than pure inadequacy. Next time, I will re{-}evaluate the dialogue to prioritize "Angry" as the dominant emotion, assigning it a higher score (e.g., 7 or 8) to reflect his defensiveness and reluctance, while lowering "Inadequate" to a secondary score (e.g., 4 or 5) to acknowledge his admitted difficulty without making it the primary driver. I will ensure "Confident" and "Relieved" remain at 0, as the text provides no evidence of these states.} \\
\\ 
\\ 
\\ 
\texttt{Your task is to predict the likely emotional responses of a character in this dialogue:} \\
\\ 
\texttt{James: I can't believe we're doing this again, Rachel. } \\
\texttt{Rachel: What, the play? Or you mean us, playing lovers?} \\
\texttt{James: Both, I guess. } \\
\texttt{Rachel: Well, it's not like we have a choice. We were cast for these roles. } \\
\texttt{James: Yeah, but it's... complicated.} \\
\texttt{Rachel: Complicated? Are you still stuck on that scene from three years ago?} \\
\texttt{James: I wouldn't say stuck. More like... haunted.} \\
\texttt{Rachel: Haunted? That's a bit dramatic, even for you, James.} \\
\texttt{James: Maybe, but it's the truth.} \\
\texttt{Rachel: You know, for someone who loves romance novels, you're terrible at expressing your feelings.} \\
\texttt{James: And you're too good at it. It's intimidating.} \\
\texttt{Rachel: Intimidating? Me? You're the one who's always retreating into your shell.} \\
\texttt{James: It's not like I do it on purpose, Rachel. It's just... hard.} \\
\texttt{{[}End dialogue{]}} \\
\\ 
\texttt{At the end of this dialogue, James would feel...} \\
\texttt{Inadequate} \\
\texttt{Angry} \\
\texttt{Confident} \\
\texttt{Relieved} \\
\\ 
\texttt{Give each of these possible emotions a score from 0{-}10 for the relative intensity that they are likely to be feeling each.} \\
\\ 
\texttt{You must output in the following format, including headings (of course, you should give your own scores), with no additional commentary:} \\
\\ 
\texttt{Inadequate: <score>} \\
\texttt{Angry: <score>} \\
\texttt{Confident: <score>} \\
\texttt{Relieved: <score>} \\
\\ 
\\ 
\texttt{{[}End of answer{]}} \\
\\ 
\texttt{Remember: zero is a valid score, meaning they are likely not feeling that emotion. You must score at least one emotion > 0.} \\
\\ 
\texttt{Your answer:} \\
            }
        }
    }
     \\ \\

    \theutterance \stepcounter{utterance}  
    & \multicolumn{4}{p{0.6\linewidth}}{
        \cellcolor[rgb]{0.561,0.737,0.561}{
            \makecell[{{p{\linewidth}}}]{
                \texttt{\tiny{[P1$\rangle$GM]}}
                \texttt{Inadequate: 7} \\
\texttt{Angry: 4} \\
\texttt{Confident: 0} \\
\texttt{Relieved: 0} \\
            }
        }
    }
    & & \\ \\

    \theutterance \stepcounter{utterance}  
    & & & \multicolumn{2}{p{0.3\linewidth}}{
        \cellcolor[rgb]{0.9,0.9,0.9}{
            \makecell[{{p{\linewidth}}}]{
                \texttt{\tiny{[GM$|$GM]}}
                \texttt{Inadequate: 7} \\
\texttt{Angry: 4} \\
\texttt{Confident: 0} \\
\texttt{Relieved: 0} \\
            }
        }
    }
    & & \\ \\

    \theutterance \stepcounter{utterance}  
    & & & \multicolumn{2}{p{0.3\linewidth}}{
        \cellcolor[rgb]{0.9,0.9,0.9}{
            \makecell[{{p{\linewidth}}}]{
                \texttt{\tiny{[GM$|$GM]}}
                \texttt{Inadequate: 8} \\
\texttt{Angry: 0} \\
\texttt{Confident: 0} \\
\texttt{Relieved: 0} \\
            }
        }
    }
    & & \\ \\

    \theutterance \stepcounter{utterance}  
    & & & \multicolumn{2}{p{0.3\linewidth}}{
        \cellcolor[rgb]{0.9,0.9,0.9}{
            \makecell[{{p{\linewidth}}}]{
                \texttt{\tiny{[GM$|$GM]}}
                \texttt{game\_result = LOSE} \\
            }
        }
    }
    & & \\ \\

    \theutterance \stepcounter{utterance}  
    & & \multicolumn{4}{p{0.6\linewidth}}{
        \cellcolor[rgb]{0.95,0.95,0.95}{
            \makecell[{{p{\linewidth}}}]{
                \texttt{\tiny{[P1$\leftrightarrow$P]}}
                \texttt{Post{-}episode Reflexion (not sent to GM):} \\
\texttt{New plan: My previous attempts failed because I misinterpreted James's "intimidation" and "retreating" as primarily stemming from a lack of self{-}worth (Inadequate), whereas the dialogue's context of a "haunted" past and defensive deflection ("It's not like I do it on purpose") points more strongly to suppressed frustration and resentment. The prompt asks for the *likely* emotional response, and James's tone is defensive and agitated rather than purely self{-}deprecating. Next time, I will re{-}evaluate the dialogue to prioritize "Angry" as the dominant emotion, assigning it the highest score (e.g., 8) to reflect his defensive stance and the "haunted" nature of his feelings. I will lower "Inadequate" to a secondary score (e.g., 3) to acknowledge his admitted difficulty without letting it overshadow the active frustration. "Confident" and "Relieved" will remain at 0. This shift from a passive self{-}doubt interpretation to an active frustration interpretation should align better with the textual evidence of his defensiveness.} \\
            }
        }
    }
    & \\ \\

\end{supertabular}
}

\end{document}
